\section{Постановка задачи}

Для данной таблицы значений функции построить методом наименьших квадратов
линейную и квадратичную полиномиальные модели, вычислить наилучшие
среднеквадратические отклонения.

Функция $y = f(x)$ задана таблицей~\ref{tab:source}. Требуется построить её
аналитическое приближение. Для этого по данным таблицы находится функция
$y = g(x)$, аппроксимирующая исходную, т.е. удовлетворяющая условию
$g(x_i) \approx y_i$, $i = 1, \ldots, n$.

Аппроксимирующая функция строится из условия минимума среднеквадратического
отклонения:
\[
\delta_2(g) = \sqrt{\frac{1}{n}\sum_{i=1}^{n}(y_i - g(x_i))^2} \to \min.
\]

В данной работе аппроксимация проводится в классе алгебраических полиномов
степени $m$:
\[
P_m(x) = a_0 + a_1 x + \cdots + a_m x^m.
\]
Коэффициенты $a_0, a_1, \ldots, a_m$ находятся из нормальной системы МНК:
\[
\sum_{j=0}^{m} a_j \sum_{i=1}^{n} x_i^{j+k} = \sum_{i=1}^{n} y_i x_i^k,
\quad k = 0, \ldots, m.
\]

\begin{table}[h!]
\centering
\begin{tabular}{|c|c|c|c|c|c|c|c|c|c|c|}
\hline
$x_i$ & 0,1 & 0,2 & 0,3 & 0,4 & 0,5 & 0,6 & 0,7 & 0,8 & 0,9 & 1,0 \\
\hline
$y_i$ & 3,15 & 3,04 & 3,02 & 2,97 & 2,87 & 2,98 & 2,81 & 2,70 & 2,66 & 2,50 \\
\hline
\hline
$x_i$ & 1,1 & 1,2 & 1,3 & 1,4 & 1,5 & 1,6 & 1,7 & 1,8 & 1,9 & 2,0 \\
\hline
$y_i$ & 2,60 & 2,36 & 2,09 & 2,07 & 2,01 & 1,81 & 1,53 & 1,64 & 1,29 & 1,11 \\
\hline
\end{tabular}
\caption{Исходные данные функции $y = f(x)$.}
\label{tab:source}
\end{table}